\todo{Maybe say some where that in this evaluation the trustworthiness property that we target the accuracy.}

\todo{Make clear that we can't validate that our methodology will work for any property but the design and results for accuracy property are already a good starting point}





\subsubsection*{Experiment 1 -- Breast Cancer Classification}
The results for this experiment are summarized in \cref{tab:cancer_results}. When both features and labels are clean, the model steadily improves and achieves high accuracy (98.68\% training, 96.49\% test), and \ac{PaTAS-TP} produces the highest trust mass across all conditions (0.88 for $\epsilon=0.1$). Clean features with corrupted labels cause test accuracy to collapse to 2.63\%, even while training accuracy remains artificially high at 99.12\%; the model learns the wrong mapping but does so confidently. Correspondingly, trust mass falls to zero regardless of $\epsilon$, correctly signaling the unreliability of the learned parameters. With vacuous labels, performance degrades moderately (65.71\% train, 74.56\% test) and trust mass remains low (0.29). Corrupted features eliminate the ability to learn in any combination: trust mass is zero in all distrust-$X$ rows and test accuracy never exceeds 43.86\%.

We evaluated the system for two values of $\epsilon$ ($0.1$ and $0.01$).
\begin{itemize}
\item For $\epsilon = 0.01$: the system is more sensitive to gradient magnitude.
      Fully trusted $X$ with trusted $Y$ yields a trust mass of 0.67, while
      uncertain $Y$ reduces it to 0.13. When $X$ is vacuous, trust stabilises
      around 0.27 for trusted $Y$ and 0.18 for uncertain $Y$.
\item For $\epsilon = 0.1$: the \todo bigger threshold makes the system less reactive
      to small gradients. Fully trusted $X$ with trusted $Y$ reaches 0.88; with
      uncertain $Y$, trust falls to 0.29. Vacuous $X$ with trusted $Y$ yields
      0.34, and with uncertain $Y$ yields 0.30. In all cases, distrust in either
      $X$ or $Y$ drives trust mass to zero.
\end{itemize}
Overall, there is a strong positive correlation between trust mass and test
accuracy. A smaller $\epsilon$ makes label trust more influential: the gap between
the trusted-$Y$ and uncertain-$Y$ rows widens as $\epsilon$ decreases.

\todo add a frame 
Two findings are worth emphasising. First, \emph{label trust matters more than
feature trust}: setting $X$ as vacuous while keeping $Y$ trusted yields a higher
trust mass (0.34) than the reverse (0.29 for trusted $X$, uncertain $Y$). Second,
\emph{distrust is more damaging than uncertainty}: replacing a fully uncertain
opinion $(0,0,1)$ for both $X$ and $Y$ with a mixed assessment $(0.25,0.25,0.5)$
reduces the trust mass from 0.30 to 0.11, showing that even partial distrust
degrades trust far more than uncertainty alone.

\todo one of the result was not commented => 0.25 0 0.75 

\todo where is Experiment 2 

\subsubsection*{Experiment 3 -- Poisoned MNIST Dataset}
As in the previous experiment, \ac{PaTAS-TP} produces ten output trust opinions.
In \cref{tab:mnistpois} we report the trust mass for label 3 (a clean class) and
label 6 (a poisoned class), alongside overall train/test accuracy and per-class
accuracy on both clean and patched test samples.

Across all patch sizes, the trust mass for label 3 consistently exceeds that for
label 6 (e.g.\ 0.836 vs.\ 0.440 for the $1\times1$ patch), confirming that
\ac{PaTAS-TP} reliably distinguishes clean from poisoned inference paths even when
overall accuracy remains high (97.66\%). As the patch grows, both trust values
decrease: the $4\times4$ and $10\times10$ patches progressively corrupt more
parameters along the inference path, dragging the trust for label 3 down from
0.836 to 0.815. The $27\times27$ patch dominates nearly the entire input, yet
trust does not collapse as dramatically as might be expected (0.651 for label 3,
0.386 for label 6), reflecting that the label-trust signal during training still
provides some protective evidence for the parameters involved in clean-class
prediction. The clear separation between the two classes is maintained across all
moderate patch sizes.\todo how does this separation evolve.

Regarding classification accuracy, patched digit 6 samples are almost entirely
misclassified regardless of patch size (0.63 to 0.84\%), confirming the backdoor
attack is effective. Patched digit 3 accuracy degrades more gradually: it drops
sharply from 83.76\% at $1\times1$ to roughly 16 to 20\% for larger patches,
reflecting the spatial overlap between the patch and discriminative features
for digit 3.



To evaluate the \ac{IPTA}, \cref{tab:mnistpoisipta} reports accuracy and trust
opinions for clean and patched samples, obtained by feedforwarding a fully trusted
(or selectively distrusted) opinion through the \ac{IPTA} instantiated from the
respective inference path.

Clean digit 3 achieves the strongest trust profile (0.811 trust, 0.189
uncertainty), consistent with high accuracy (98.02\%). Clean digit 6 already
exhibits markedly higher uncertainty (0.570) despite comparable accuracy (97.70\%),
reflecting that the parameters on the digit-6 inference path were exposed to
poisoned training samples and are therefore less trusted.

For patched samples, the degradation is clear. Patched digit 6 accuracy collapses
to 0.84\%, and trust drops slightly to 0.408, with uncertainty dominating at 0.592.
Notably, when patch pixels are explicitly distrusted in the input $T_x$, a small
but non-zero distrust component (0.008) emerges in the output, providing an
actionable warning signal unavailable from accuracy alone. The effect is more
pronounced for patched digit 3: accuracy falls to 19.80\% yet trust remains
relatively high at 0.738 (trusted $T_x$) and 0.724 (patch distrusted), reflecting
that the digit-3 inference path passes through less poisoned parameters. Even so,
explicitly distrusting the patch pixels introduces a distrust mass of 0.014 and
reduces the trust mass, showing that \ac{IPTA} propagates localised input
scepticism into the output opinion.

\todo the more distrust in patch 3, as we can see is due to the fact that the distrust in the patch translate a part of the output trust mass into a distrust mass. And since trusted patch 3 as a bigger trust mass, the distrusted one naturally ends with bigger distrust mass.
\todo maybe just remove this part as it doesn't give very useful insight 



\todo{add a frame, same for all results}

These results demonstrate that \ac{PaTAS-TP} not only correlates with accuracy
but provides complementary information: it identifies which inference paths are
structurally less reliable \todo (clean label vs.\ non clean, and pacthed image vs non patched one) and responds to targeted input
distrust in a semantically meaningful way.






% Replace the stale IPTA reference:
Results show a degradation in trust for adversarial samples: benign digit 3
achieved $t=0.811, d=0, u=0.189$, while patched digit 3 dropped to
$t=0.738, d=0, u=0.262$ (trusted $T_x$), or $t=0.724, d=0.014, u=0.262$ \todo 
when patch pixels are explicitly distrusted.

% Replace the stale breast cancer reference:
In the breast cancer task (\cref{tab:cancer_results}), assigning a mixed trust
profile $(0.25, 0, 0.75)$ to features with trusted labels yields a final trust
mass of $0.38$ for $\epsilon=0.1$, while the vacuous-features/trusted-labels
case achieves $0.34$ with comparable test accuracy (95.61\%). More strikingly,
the fully trusted scenario (trust mass $0.88$) coincides with the highest test
accuracy (96.49\%), whereas the mixed distrust scenario $(0.25, 0.25, 0.5)$
reduces trust mass to $0.11$ while test accuracy remains at 88.60\%, confirming
that partial distrust erodes trust assessments far more than uncertainty alone.
\todo last par to be improved 
